\scanpage{34}
% L1–5
(i) signifie que $D(A)\simeq H^n(A)$, i.e. (ii).

ii) $\Rightarrow$ (iii), d'autre part (iii) $\Rightarrow H^n(A)\simeq D(A)^m$ i.e. $A^m$ est \deleted{dualisant} fondamental. En vertu de prop. 9, on a $\lg_{\mathfrak p}(A^m)=\lg_{\mathfrak p}(A)$, d'où $m=1$, d'où (iii) $\Rightarrow$ (i). \deleted{Il reste à \unsure{prouver}}

% L5–8
Evident (i) $\Rightarrow$ (iv) par prop. 9, réciproque : $H^n(A)$ \unsure{sera|s'in.|comon.} $\subset D(A)$, l'annulateur de $H^n(A)$ est nul d'après l'hyp., donc $H^n(A)=D(A)$.

% L9–10
\deleted{\uline{Définition 2} Un tel $A$ est dit anneau de Gorenstein. S'il est de plus Cohen-Macaulay}

% L11–15, crossed draft
\deleted{\uline{Th. 5} Un anneau de Gorenstein est Coh. Mac}

\deleted{En effet, il est équidim, \unsure{donc} si \unsure{sa dim $n$}} \deleted{Réc. sur $n=\dim A$, trivial si $n=0$. Si $n>0$, $A$ étant équidim, il y a un $x\in A$ \unsure{non|t non} diviseur de $0$. \unsure{Considérons} Alors \deleted{la suite exacte} \inserted{la suite spectrale}}

% L16–20
\deleted{(où $N$ est un \deleted{\unsure{$A$}} \inserted{$A/x$}-$A/xA$-\unsure{module}) \unsure{de} \unsure{donne}}
\[
 \deleted{0\longrightarrow H^{n}(}\qquad \deleted{H^{n-1}(\,)}
\]
\[
 \deleted{\operatorname{Ext}_{A}^{*}(N,A\deleted{\unsure{/x}})\Longleftarrow
 \operatorname{Ext}_{A/x}^{*}(N,\operatorname{Ext}_{A}^{*}(A/x,A))}
\]
\deleted{\unsure{donne}, puisque} d'après le cor. au lemme 6,

% L21–25
\uline{Prop 16} bis \inserted{$+\,A$ est Cohen-Macaulay}, ces conditions équivalent aussi aux suivantes

\deleted{\unsure{b')|a')} \unsure{(ii)bis} Le foncteur $N\longrightarrow\operatorname{Ext}_{A}^{n}(N,A)$ est dualisant} \deleted{\unsure{ou que}} \deleted{[ou simplement exact]} \inserted{\deleted{exact}}

\inserted{c')} \deleted{\unsure{d)}} $\operatorname{Ext}^{n}(k,A)\simeq k$

\deleted{\unsure{e)}} $\operatorname{Ext}^{n+1}(k,A)=0$

% archive sheet number: 33
